\section{西夏的河西社会与文书国家}

\subsection{户籍与地方分类}

982年李继迁脱离羁縻。这一账本锚点把。国家记录通常从诏令、战争、改制、灾异或官员活动进入，地方社会却在同一节点中面对登记、运输、价格、婚姻、服役、避险和迁居等更缓慢的选择。因而，本节不将制度名称当作已经实现的秩序，而是追问它需要哪些书写、人员、道路和资源才能抵达具体地点。

本段以。编年、本纪和制度材料适合确认政治年序与规范语言；契约、讼牍、墓志、碑刻、寺院档案、钱币、城址、窑址、河道与港口遗存则分别保存有限的地方证据。它们的差异不是缺陷而是方法前提：一条律令说明应当如何分类，一份案件才显示某次争执被如何说出；一座遗址说明被发掘范围内的生产或居住，不能直接等于全城或全境。

阶层关系在材料中从不以完整名单出现。官府能够登记户、役、田、商税或军人，却未必记录同一家庭中妇女、儿童、雇佣者、佃户、债务人与流动亲属的实际分工。精英家族留下较多墓志和文集，也使他们的语言容易被误作社会常态。写作时须把国家分类、家族自述和地方实践分开：它们可以重叠，却不能被一概视为同一身份。

环境与区域差异进一步改变制度的效果。河流、山地、草场、港口、旱涝、疫病和交通季节性会影响征粮、赴役、贸易与逃亡；但灾害记录带有观察地点和政治修辞，不能由一条灾异推论整个政权的经济或人口。把环境证据同道路、仓储、市场和家庭策略并置，能够说明为何相同的命令在都城附近、边寨、河谷和海域产生不同后果。

文化与宗教亦属于这一社会基础。经卷、题记、书院记录、学校条目、法律文本和技术书显示书写与知识如何被保管、传递或用于行政；它们不证明所有居民都读过同一文本，也不能将捐施或立碑的愿望直接等同于普遍信仰。多语文字、官号和族名尤其必须保留原始语境，不能把一次书写实践简化为固定的族群本质。

因此，，而在于令：谁有资格登记和裁决，谁承担风险，谁可借助市场、亲属、寺院或地方网络重新组织生活。材料未能覆盖之处，本文明确保留限度，不以流畅叙事取代证据。

\subsection{土地、劳役与家庭经济}

1009年经营夏州灵州。这一账本锚点把。国家记录通常从诏令、战争、改制、灾异或官员活动进入，地方社会却在同一节点中面对登记、运输、价格、婚姻、服役、避险和迁居等更缓慢的选择。因而，本节不将制度名称当作已经实现的秩序，而是追问它需要哪些书写、人员、道路和资源才能抵达具体地点。

本段以。编年、本纪和制度材料适合确认政治年序与规范语言；契约、讼牍、墓志、碑刻、寺院档案、钱币、城址、窑址、河道与港口遗存则分别保存有限的地方证据。它们的差异不是缺陷而是方法前提：一条律令说明应当如何分类，一份案件才显示某次争执被如何说出；一座遗址说明被发掘范围内的生产或居住，不能直接等于全城或全境。

阶层关系在材料中从不以完整名单出现。官府能够登记户、役、田、商税或军人，却未必记录同一家庭中妇女、儿童、雇佣者、佃户、债务人与流动亲属的实际分工。精英家族留下较多墓志和文集，也使他们的语言容易被误作社会常态。写作时须把国家分类、家族自述和地方实践分开：它们可以重叠，却不能被一概视为同一身份。

环境与区域差异进一步改变制度的效果。河流、山地、草场、港口、旱涝、疫病和交通季节性会影响征粮、赴役、贸易与逃亡；但灾害记录带有观察地点和政治修辞，不能由一条灾异推论整个政权的经济或人口。把环境证据同道路、仓储、市场和家庭策略并置，能够说明为何相同的命令在都城附近、边寨、河谷和海域产生不同后果。

文化与宗教亦属于这一社会基础。经卷、题记、书院记录、学校条目、法律文本和技术书显示书写与知识如何被保管、传递或用于行政；它们不证明所有居民都读过同一文本，也不能将捐施或立碑的愿望直接等同于普遍信仰。多语文字、官号和族名尤其必须保留原始语境，不能把一次书写实践简化为固定的族群本质。

因此，，而在于令：谁有资格登记和裁决，谁承担风险，谁可借助市场、亲属、寺院或地方网络重新组织生活。材料未能覆盖之处，本文明确保留限度，不以流畅叙事取代证据。

\subsection{城市、市场与工匠}

1034年改国号年号。这一账本锚点把。国家记录通常从诏令、战争、改制、灾异或官员活动进入，地方社会却在同一节点中面对登记、运输、价格、婚姻、服役、避险和迁居等更缓慢的选择。因而，本节不将制度名称当作已经实现的秩序，而是追问它需要哪些书写、人员、道路和资源才能抵达具体地点。

本段以。编年、本纪和制度材料适合确认政治年序与规范语言；契约、讼牍、墓志、碑刻、寺院档案、钱币、城址、窑址、河道与港口遗存则分别保存有限的地方证据。它们的差异不是缺陷而是方法前提：一条律令说明应当如何分类，一份案件才显示某次争执被如何说出；一座遗址说明被发掘范围内的生产或居住，不能直接等于全城或全境。

阶层关系在材料中从不以完整名单出现。官府能够登记户、役、田、商税或军人，却未必记录同一家庭中妇女、儿童、雇佣者、佃户、债务人与流动亲属的实际分工。精英家族留下较多墓志和文集，也使他们的语言容易被误作社会常态。写作时须把国家分类、家族自述和地方实践分开：它们可以重叠，却不能被一概视为同一身份。

环境与区域差异进一步改变制度的效果。河流、山地、草场、港口、旱涝、疫病和交通季节性会影响征粮、赴役、贸易与逃亡；但灾害记录带有观察地点和政治修辞，不能由一条灾异推论整个政权的经济或人口。把环境证据同道路、仓储、市场和家庭策略并置，能够说明为何相同的命令在都城附近、边寨、河谷和海域产生不同后果。

文化与宗教亦属于这一社会基础。经卷、题记、书院记录、学校条目、法律文本和技术书显示书写与知识如何被保管、传递或用于行政；它们不证明所有居民都读过同一文本，也不能将捐施或立碑的愿望直接等同于普遍信仰。多语文字、官号和族名尤其必须保留原始语境，不能把一次书写实践简化为固定的族群本质。

因此，，而在于令：谁有资格登记和裁决，谁承担风险，谁可借助市场、亲属、寺院或地方网络重新组织生活。材料未能覆盖之处，本文明确保留限度，不以流畅叙事取代证据。

\subsection{水利、灾害与生计}

1036年扩张甘凉。这一账本锚点把。国家记录通常从诏令、战争、改制、灾异或官员活动进入，地方社会却在同一节点中面对登记、运输、价格、婚姻、服役、避险和迁居等更缓慢的选择。因而，本节不将制度名称当作已经实现的秩序，而是追问它需要哪些书写、人员、道路和资源才能抵达具体地点。

本段以。编年、本纪和制度材料适合确认政治年序与规范语言；契约、讼牍、墓志、碑刻、寺院档案、钱币、城址、窑址、河道与港口遗存则分别保存有限的地方证据。它们的差异不是缺陷而是方法前提：一条律令说明应当如何分类，一份案件才显示某次争执被如何说出；一座遗址说明被发掘范围内的生产或居住，不能直接等于全城或全境。

阶层关系在材料中从不以完整名单出现。官府能够登记户、役、田、商税或军人，却未必记录同一家庭中妇女、儿童、雇佣者、佃户、债务人与流动亲属的实际分工。精英家族留下较多墓志和文集，也使他们的语言容易被误作社会常态。写作时须把国家分类、家族自述和地方实践分开：它们可以重叠，却不能被一概视为同一身份。

环境与区域差异进一步改变制度的效果。河流、山地、草场、港口、旱涝、疫病和交通季节性会影响征粮、赴役、贸易与逃亡；但灾害记录带有观察地点和政治修辞，不能由一条灾异推论整个政权的经济或人口。把环境证据同道路、仓储、市场和家庭策略并置，能够说明为何相同的命令在都城附近、边寨、河谷和海域产生不同后果。

文化与宗教亦属于这一社会基础。经卷、题记、书院记录、学校条目、法律文本和技术书显示书写与知识如何被保管、传递或用于行政；它们不证明所有居民都读过同一文本，也不能将捐施或立碑的愿望直接等同于普遍信仰。多语文字、官号和族名尤其必须保留原始语境，不能把一次书写实践简化为固定的族群本质。

因此，，而在于令：谁有资格登记和裁决，谁承担风险，谁可借助市场、亲属、寺院或地方网络重新组织生活。材料未能覆盖之处，本文明确保留限度，不以流畅叙事取代证据。

\subsection{边防、道路与流动}

1038年西夏建国。这一账本锚点把。国家记录通常从诏令、战争、改制、灾异或官员活动进入，地方社会却在同一节点中面对登记、运输、价格、婚姻、服役、避险和迁居等更缓慢的选择。因而，本节不将制度名称当作已经实现的秩序，而是追问它需要哪些书写、人员、道路和资源才能抵达具体地点。

本段以。编年、本纪和制度材料适合确认政治年序与规范语言；契约、讼牍、墓志、碑刻、寺院档案、钱币、城址、窑址、河道与港口遗存则分别保存有限的地方证据。它们的差异不是缺陷而是方法前提：一条律令说明应当如何分类，一份案件才显示某次争执被如何说出；一座遗址说明被发掘范围内的生产或居住，不能直接等于全城或全境。

阶层关系在材料中从不以完整名单出现。官府能够登记户、役、田、商税或军人，却未必记录同一家庭中妇女、儿童、雇佣者、佃户、债务人与流动亲属的实际分工。精英家族留下较多墓志和文集，也使他们的语言容易被误作社会常态。写作时须把国家分类、家族自述和地方实践分开：它们可以重叠，却不能被一概视为同一身份。

环境与区域差异进一步改变制度的效果。河流、山地、草场、港口、旱涝、疫病和交通季节性会影响征粮、赴役、贸易与逃亡；但灾害记录带有观察地点和政治修辞，不能由一条灾异推论整个政权的经济或人口。把环境证据同道路、仓储、市场和家庭策略并置，能够说明为何相同的命令在都城附近、边寨、河谷和海域产生不同后果。

文化与宗教亦属于这一社会基础。经卷、题记、书院记录、学校条目、法律文本和技术书显示书写与知识如何被保管、传递或用于行政；它们不证明所有居民都读过同一文本，也不能将捐施或立碑的愿望直接等同于普遍信仰。多语文字、官号和族名尤其必须保留原始语境，不能把一次书写实践简化为固定的族群本质。

因此，，而在于令：谁有资格登记和裁决，谁承担风险，谁可借助市场、亲属、寺院或地方网络重新组织生活。材料未能覆盖之处，本文明确保留限度，不以流畅叙事取代证据。

\subsection{寺院、教育与知识}

1040--1044年宋夏战争和议。这一账本锚点把。国家记录通常从诏令、战争、改制、灾异或官员活动进入，地方社会却在同一节点中面对登记、运输、价格、婚姻、服役、避险和迁居等更缓慢的选择。因而，本节不将制度名称当作已经实现的秩序，而是追问它需要哪些书写、人员、道路和资源才能抵达具体地点。

本段以。编年、本纪和制度材料适合确认政治年序与规范语言；契约、讼牍、墓志、碑刻、寺院档案、钱币、城址、窑址、河道与港口遗存则分别保存有限的地方证据。它们的差异不是缺陷而是方法前提：一条律令说明应当如何分类，一份案件才显示某次争执被如何说出；一座遗址说明被发掘范围内的生产或居住，不能直接等于全城或全境。

阶层关系在材料中从不以完整名单出现。官府能够登记户、役、田、商税或军人，却未必记录同一家庭中妇女、儿童、雇佣者、佃户、债务人与流动亲属的实际分工。精英家族留下较多墓志和文集，也使他们的语言容易被误作社会常态。写作时须把国家分类、家族自述和地方实践分开：它们可以重叠，却不能被一概视为同一身份。

环境与区域差异进一步改变制度的效果。河流、山地、草场、港口、旱涝、疫病和交通季节性会影响征粮、赴役、贸易与逃亡；但灾害记录带有观察地点和政治修辞，不能由一条灾异推论整个政权的经济或人口。把环境证据同道路、仓储、市场和家庭策略并置，能够说明为何相同的命令在都城附近、边寨、河谷和海域产生不同后果。

文化与宗教亦属于这一社会基础。经卷、题记、书院记录、学校条目、法律文本和技术书显示书写与知识如何被保管、传递或用于行政；它们不证明所有居民都读过同一文本，也不能将捐施或立碑的愿望直接等同于普遍信仰。多语文字、官号和族名尤其必须保留原始语境，不能把一次书写实践简化为固定的族群本质。

因此，，而在于令：谁有资格登记和裁决，谁承担风险，谁可借助市场、亲属、寺院或地方网络重新组织生活。材料未能覆盖之处，本文明确保留限度，不以流畅叙事取代证据。

\subsection{司法、文书与基层秩序}

1048年元昊遇害。这一账本锚点把。国家记录通常从诏令、战争、改制、灾异或官员活动进入，地方社会却在同一节点中面对登记、运输、价格、婚姻、服役、避险和迁居等更缓慢的选择。因而，本节不将制度名称当作已经实现的秩序，而是追问它需要哪些书写、人员、道路和资源才能抵达具体地点。

本段以。编年、本纪和制度材料适合确认政治年序与规范语言；契约、讼牍、墓志、碑刻、寺院档案、钱币、城址、窑址、河道与港口遗存则分别保存有限的地方证据。它们的差异不是缺陷而是方法前提：一条律令说明应当如何分类，一份案件才显示某次争执被如何说出；一座遗址说明被发掘范围内的生产或居住，不能直接等于全城或全境。

阶层关系在材料中从不以完整名单出现。官府能够登记户、役、田、商税或军人，却未必记录同一家庭中妇女、儿童、雇佣者、佃户、债务人与流动亲属的实际分工。精英家族留下较多墓志和文集，也使他们的语言容易被误作社会常态。写作时须把国家分类、家族自述和地方实践分开：它们可以重叠，却不能被一概视为同一身份。

环境与区域差异进一步改变制度的效果。河流、山地、草场、港口、旱涝、疫病和交通季节性会影响征粮、赴役、贸易与逃亡；但灾害记录带有观察地点和政治修辞，不能由一条灾异推论整个政权的经济或人口。把环境证据同道路、仓储、市场和家庭策略并置，能够说明为何相同的命令在都城附近、边寨、河谷和海域产生不同后果。

文化与宗教亦属于这一社会基础。经卷、题记、书院记录、学校条目、法律文本和技术书显示书写与知识如何被保管、传递或用于行政；它们不证明所有居民都读过同一文本，也不能将捐施或立碑的愿望直接等同于普遍信仰。多语文字、官号和族名尤其必须保留原始语境，不能把一次书写实践简化为固定的族群本质。

因此，，而在于令：谁有资格登记和裁决，谁承担风险，谁可借助市场、亲属、寺院或地方网络重新组织生活。材料未能覆盖之处，本文明确保留限度，不以流畅叙事取代证据。

\subsection{性别、婚姻与继承}

1082年永乐城攻守。这一账本锚点把。国家记录通常从诏令、战争、改制、灾异或官员活动进入，地方社会却在同一节点中面对登记、运输、价格、婚姻、服役、避险和迁居等更缓慢的选择。因而，本节不将制度名称当作已经实现的秩序，而是追问它需要哪些书写、人员、道路和资源才能抵达具体地点。

本段以。编年、本纪和制度材料适合确认政治年序与规范语言；契约、讼牍、墓志、碑刻、寺院档案、钱币、城址、窑址、河道与港口遗存则分别保存有限的地方证据。它们的差异不是缺陷而是方法前提：一条律令说明应当如何分类，一份案件才显示某次争执被如何说出；一座遗址说明被发掘范围内的生产或居住，不能直接等于全城或全境。

阶层关系在材料中从不以完整名单出现。官府能够登记户、役、田、商税或军人，却未必记录同一家庭中妇女、儿童、雇佣者、佃户、债务人与流动亲属的实际分工。精英家族留下较多墓志和文集，也使他们的语言容易被误作社会常态。写作时须把国家分类、家族自述和地方实践分开：它们可以重叠，却不能被一概视为同一身份。

环境与区域差异进一步改变制度的效果。河流、山地、草场、港口、旱涝、疫病和交通季节性会影响征粮、赴役、贸易与逃亡；但灾害记录带有观察地点和政治修辞，不能由一条灾异推论整个政权的经济或人口。把环境证据同道路、仓储、市场和家庭策略并置，能够说明为何相同的命令在都城附近、边寨、河谷和海域产生不同后果。

文化与宗教亦属于这一社会基础。经卷、题记、书院记录、学校条目、法律文本和技术书显示书写与知识如何被保管、传递或用于行政；它们不证明所有居民都读过同一文本，也不能将捐施或立碑的愿望直接等同于普遍信仰。多语文字、官号和族名尤其必须保留原始语境，不能把一次书写实践简化为固定的族群本质。

因此，，而在于令：谁有资格登记和裁决，谁承担风险，谁可借助市场、亲属、寺院或地方网络重新组织生活。材料未能覆盖之处，本文明确保留限度，不以流畅叙事取代证据。

\subsection{区域交换与外来者}

1144年金夏边境安排。这一账本锚点把。国家记录通常从诏令、战争、改制、灾异或官员活动进入，地方社会却在同一节点中面对登记、运输、价格、婚姻、服役、避险和迁居等更缓慢的选择。因而，本节不将制度名称当作已经实现的秩序，而是追问它需要哪些书写、人员、道路和资源才能抵达具体地点。

本段以。编年、本纪和制度材料适合确认政治年序与规范语言；契约、讼牍、墓志、碑刻、寺院档案、钱币、城址、窑址、河道与港口遗存则分别保存有限的地方证据。它们的差异不是缺陷而是方法前提：一条律令说明应当如何分类，一份案件才显示某次争执被如何说出；一座遗址说明被发掘范围内的生产或居住，不能直接等于全城或全境。

阶层关系在材料中从不以完整名单出现。官府能够登记户、役、田、商税或军人，却未必记录同一家庭中妇女、儿童、雇佣者、佃户、债务人与流动亲属的实际分工。精英家族留下较多墓志和文集，也使他们的语言容易被误作社会常态。写作时须把国家分类、家族自述和地方实践分开：它们可以重叠，却不能被一概视为同一身份。

环境与区域差异进一步改变制度的效果。河流、山地、草场、港口、旱涝、疫病和交通季节性会影响征粮、赴役、贸易与逃亡；但灾害记录带有观察地点和政治修辞，不能由一条灾异推论整个政权的经济或人口。把环境证据同道路、仓储、市场和家庭策略并置，能够说明为何相同的命令在都城附近、边寨、河谷和海域产生不同后果。

文化与宗教亦属于这一社会基础。经卷、题记、书院记录、学校条目、法律文本和技术书显示书写与知识如何被保管、传递或用于行政；它们不证明所有居民都读过同一文本，也不能将捐施或立碑的愿望直接等同于普遍信仰。多语文字、官号和族名尤其必须保留原始语境，不能把一次书写实践简化为固定的族群本质。

因此，，而在于令：谁有资格登记和裁决，谁承担风险，谁可借助市场、亲属、寺院或地方网络重新组织生活。材料未能覆盖之处，本文明确保留限度，不以流畅叙事取代证据。

\subsection{迁徙、安置与社会记忆}

1209--1227年蒙古压力。这一账本锚点把。国家记录通常从诏令、战争、改制、灾异或官员活动进入，地方社会却在同一节点中面对登记、运输、价格、婚姻、服役、避险和迁居等更缓慢的选择。因而，本节不将制度名称当作已经实现的秩序，而是追问它需要哪些书写、人员、道路和资源才能抵达具体地点。

本段以。编年、本纪和制度材料适合确认政治年序与规范语言；契约、讼牍、墓志、碑刻、寺院档案、钱币、城址、窑址、河道与港口遗存则分别保存有限的地方证据。它们的差异不是缺陷而是方法前提：一条律令说明应当如何分类，一份案件才显示某次争执被如何说出；一座遗址说明被发掘范围内的生产或居住，不能直接等于全城或全境。

阶层关系在材料中从不以完整名单出现。官府能够登记户、役、田、商税或军人，却未必记录同一家庭中妇女、儿童、雇佣者、佃户、债务人与流动亲属的实际分工。精英家族留下较多墓志和文集，也使他们的语言容易被误作社会常态。写作时须把国家分类、家族自述和地方实践分开：它们可以重叠，却不能被一概视为同一身份。

环境与区域差异进一步改变制度的效果。河流、山地、草场、港口、旱涝、疫病和交通季节性会影响征粮、赴役、贸易与逃亡；但灾害记录带有观察地点和政治修辞，不能由一条灾异推论整个政权的经济或人口。把环境证据同道路、仓储、市场和家庭策略并置，能够说明为何相同的命令在都城附近、边寨、河谷和海域产生不同后果。

文化与宗教亦属于这一社会基础。经卷、题记、书院记录、学校条目、法律文本和技术书显示书写与知识如何被保管、传递或用于行政；它们不证明所有居民都读过同一文本，也不能将捐施或立碑的愿望直接等同于普遍信仰。多语文字、官号和族名尤其必须保留原始语境，不能把一次书写实践简化为固定的族群本质。

因此，，而在于令：谁有资格登记和裁决，谁承担风险，谁可借助市场、亲属、寺院或地方网络重新组织生活。材料未能覆盖之处，本文明确保留限度，不以流畅叙事取代证据。

\subsection{国家尺度与地方尺度之间}

国家制度之所以看似完整，常因史书和政书把不同年份、不同地域的条目整理为稳定名目。实际生活中的征收、裁判、教育、宗教和市场却由人员往返、文书抄写、道路维护和地方关系维系。某项制度被反复申明，可能意味着官府在追求一致性，也可能暴露其难以落实。将复申直接写为成功或失败，都忽略了地方机构如何在有限资源中选择性执行。

地方行动者并非总是对抗国家。家族、寺院、行会、学校、军镇和乡里可提供登记、救济、借贷、调解与知识传播的渠道，也可能成为催科、征发和纪律进入社区的接口。它们既不能被浪漫化为自治社会，也不应被简化为中枢命令的被动工具。每一份材料的形成者、对象和保存过程，决定了我们能够看见何种合作、何种冲突以及何种沉默。

比较各区域时，应让时间差保持可见。都城的制度创新可能多年后才在边地出现，战争前线的流民和军需也可能先于法令改变市场。灾害后的复耕、重建和迁居各有节奏，不能因一个年份的政令就假定社会同步恢复。只有在这样的慢时间中，阶层、性别、职业与地域才不至于变成脱离材料的标签。

\subsection{证据的保存与失落}

档案、碑刻、墓志、契约和遗址并不是自然留下的窗口。它们经过官府筛选、家族保存、寺院收藏、后世抄录、盗掘或考古发掘，才进入今天的研究。保存过程造成的偏向要求我们既使用细节，也不把细节外推为普遍规律。对没有留下文字的人群，物质遗存与环境证据可提供间接约束，但不能替他们编造统一的声音。

制度史与社会史的结合，不是把一条法令附上几件个案，而是询问两者能否在同一地点、同一时间、同一类行动者身上相遇。若不能相遇，差距本身就是结论：国家的分类可能尚未抵达，地方的实践也可能未被国家记录。这样的限度并不削弱叙事，反而使区域差异、家庭策略和政治选择具有可以检验的尺度。
\subsection{证据的保存与失落}

档案、碑刻、墓志、契约和遗址并不是自然留下的窗口。它们经过官府筛选、家族保存、寺院收藏、后世抄录、盗掘或考古发掘，才进入今天的研究。保存过程造成的偏向要求我们既使用细节，也不把细节外推为普遍规律。对没有留下文字的人群，物质遗存与环境证据可提供间接约束，但不能替他们编造统一的声音。

制度史与社会史的结合，不是把一条法令附上几件个案，而是询问两者能否在同一地点、同一时间、同一类行动者身上相遇。若不能相遇，差距本身就是结论：国家的分类可能尚未抵达，地方的实践也可能未被国家记录。这样的限度并不削弱叙事，反而使区域差异、家庭策略和政治选择具有可以检验的尺度。

\subsection{方法的自我限制}

历史解释必须与材料的分辨率相称。可以确定的年序不必能确定动机，可以确认的制度名称不必能确认执行规模，可以识别的遗址不必能重建全部社区。将这些层次逐一分开，既能避免把后见之明投回当时，也能使不同政权、区域和人群之间的比较更为可靠。叙事所保留的不确定性，不是空白，而是提醒读者继续追问证据来自谁、服务何种目的、又遗漏了谁。
\subsection{方法的自我限制}

历史解释必须与材料的分辨率相称。可以确定的年序不必能确定动机，可以确认的制度名称不必能确认执行规模，可以识别的遗址不必能重建全部社区。将这些层次逐一分开，既能避免把后见之明投回当时，也能使不同政权、区域和人群之间的比较更为可靠。叙事所保留的不确定性，不是空白，而是提醒读者继续追问证据来自谁、服务何种目的、又遗漏了谁。

